\documentclass[11pt, letterpaper]{article}
\usepackage[margin=1in]{geometry}
\usepackage{booktabs}
\usepackage{array}
\usepackage{tabularx}
\usepackage[table]{xcolor}
\usepackage{multirow}
\usepackage{enumitem}
\usepackage{titlesec}
\usepackage[hidelinks]{hyperref}
\usepackage{amsmath}
\usepackage{amssymb}
\usepackage[expansion=false]{microtype}
\usepackage{tcolorbox}
\usepackage{multicol}
\usepackage{titletoc}
\usepackage{graphicx}
\usepackage{listings}
\usepackage{courier}

\tcbuselibrary{skins, breakable}

% ── Colors ───────────────────────────────────────────────────────────────────
\definecolor{sectionblue}{RGB}{30, 90, 160}
\definecolor{sectiongreen}{RGB}{30, 130, 80}
\definecolor{sectionpurple}{RGB}{100, 40, 140}
\definecolor{sectionorange}{RGB}{180, 90, 20}
\definecolor{sectionteal}{RGB}{20, 120, 120}
\definecolor{sectionred}{RGB}{160, 30, 40}
\definecolor{tableheader}{RGB}{220, 230, 245}
\definecolor{tableheadergreen}{RGB}{210, 240, 220}
\definecolor{tableheaderorange}{RGB}{255, 235, 210}
\definecolor{tableheaderteal}{RGB}{210, 240, 240}
\definecolor{tableheaderred}{RGB}{250, 220, 222}
\definecolor{tableheaderpurple}{RGB}{235, 220, 245}
\definecolor{rowA}{RGB}{252, 252, 252}
\definecolor{rowB}{RGB}{237, 244, 255}
\definecolor{rowBg}{RGB}{237, 248, 241}
\definecolor{rowOr}{RGB}{255, 245, 232}
\definecolor{rowTe}{RGB}{232, 248, 248}
\definecolor{rowPu}{RGB}{245, 235, 255}
\definecolor{rowRe}{RGB}{255, 238, 238}
\definecolor{partbg}{RGB}{240, 244, 255}
\definecolor{part2bg}{RGB}{240, 252, 244}
\definecolor{part3bg}{RGB}{248, 240, 255}
\definecolor{part4bg}{RGB}{255, 245, 232}
\definecolor{part5bg}{RGB}{232, 248, 248}
\definecolor{part6bg}{RGB}{255, 240, 240}
\definecolor{notebg}{RGB}{255, 252, 228}
\definecolor{noteborder}{RGB}{180, 148, 30}
\definecolor{codebg}{RGB}{245, 245, 250}
\definecolor{codeborder}{RGB}{100, 100, 140}

% ── Section switchers ─────────────────────────────────────────────────────────
\newcommand{\bluesections}{%
  \titleformat{\section}{\large\bfseries\color{sectionblue}}{}{0em}{}[\color{sectionblue}\titlerule]%
  \titleformat{\subsection}{\normalsize\bfseries\color{sectionblue!80!black}}{}{0em}{}%
}
\newcommand{\greensections}{%
  \titleformat{\section}{\large\bfseries\color{sectiongreen}}{}{0em}{}[\color{sectiongreen}\titlerule]%
  \titleformat{\subsection}{\normalsize\bfseries\color{sectiongreen!80!black}}{}{0em}{}%
}
\newcommand{\purplesections}{%
  \titleformat{\section}{\large\bfseries\color{sectionpurple}}{}{0em}{}[\color{sectionpurple}\titlerule]%
  \titleformat{\subsection}{\normalsize\bfseries\color{sectionpurple!80!black}}{}{0em}{}%
}
\newcommand{\orangesections}{%
  \titleformat{\section}{\large\bfseries\color{sectionorange}}{}{0em}{}[\color{sectionorange}\titlerule]%
  \titleformat{\subsection}{\normalsize\bfseries\color{sectionorange!80!black}}{}{0em}{}%
}
\newcommand{\tealsections}{%
  \titleformat{\section}{\large\bfseries\color{sectionteal}}{}{0em}{}[\color{sectionteal}\titlerule]%
  \titleformat{\subsection}{\normalsize\bfseries\color{sectionteal!80!black}}{}{0em}{}%
}
\newcommand{\redsections}{%
  \titleformat{\section}{\large\bfseries\color{sectionred}}{}{0em}{}[\color{sectionred}\titlerule]%
  \titleformat{\subsection}{\normalsize\bfseries\color{sectionred!80!black}}{}{0em}{}%
}

% ── Part banner ───────────────────────────────────────────────────────────────
\newcommand{\partbanner}[3]{%
  \begin{tcolorbox}[colback=#1, colframe=#2, boxrule=1.2pt, arc=4pt,
    left=8pt, right=8pt, top=6pt, bottom=6pt]
    \begin{center}{\LARGE\bfseries #3}\end{center}
  \end{tcolorbox}%
}

% ── Code listing style ────────────────────────────────────────────────────────
\lstdefinestyle{mips}{
  basicstyle=\ttfamily\small,
  backgroundcolor=\color{codebg},
  frame=single,
  rulecolor=\color{codeborder},
  keywordstyle=\bfseries\color{sectionblue},
  commentstyle=\itshape\color{sectiongreen},
  stringstyle=\color{sectionred},
  comment=[l]{\#},
  morekeywords={add,addi,sub,and,andi,or,ori,nor,sll,srl,lw,sw,beq,bne,j,jal,jr,slt,slti,lui,mfhi,mflo,mult,div},
  numbers=left,
  numberstyle=\tiny\color{gray},
  numbersep=6pt,
  xleftmargin=16pt,
  breaklines=true,
  tabsize=4,
  showstringspaces=false
}
\lstdefinestyle{x86}{
  basicstyle=\ttfamily\small,
  backgroundcolor=\color{codebg},
  frame=single,
  rulecolor=\color{codeborder},
  keywordstyle=\bfseries\color{sectionred},
  commentstyle=\itshape\color{sectiongreen},
  stringstyle=\color{sectionpurple},
  comment=[l]{;},
  morekeywords={mov,push,pop,call,ret,add,sub,mul,imul,div,idiv,and,or,xor,not,
                cmp,test,jmp,je,jne,jl,jle,jg,jge,jz,jnz,lea,nop,syscall,
                int,inc,dec,shl,shr,sar,sal,neg,cbw,cwde,cdqe,
                rax,rbx,rcx,rdx,rsi,rdi,rsp,rbp,r8,r9,r10,r11,r12,r13,r14,r15,
                eax,ebx,ecx,edx,esi,edi,esp,ebp,
                ax,bx,cx,dx,si,di,sp,bp,
                al,bl,cl,dl,ah,bh,ch,dh,
                QWORD,DWORD,WORD,BYTE,PTR,OFFSET},
  numbers=left,
  numberstyle=\tiny\color{gray},
  numbersep=6pt,
  xleftmargin=16pt,
  breaklines=true,
  tabsize=4,
  showstringspaces=false
}

\setlength{\parskip}{4pt}
\setlength{\parindent}{0pt}
\newcommand{\divider}{\par\noindent\rule{\textwidth}{0.4pt}\par\vspace{2pt}}
\hbadness=10000
\hfuzz=5pt

% ── TOC styling ───────────────────────────────────────────────────────────────
\titlecontents{section}[1.5em]{\smallskip}
  {\contentslabel{1.5em}}{\hspace*{-1.5em}}
  {\titlerule*[6pt]{.}\contentspage}
\titlecontents{subsection}[3.5em]{}
  {\contentslabel{2em}}{\hspace*{-2em}}
  {\titlerule*[6pt]{.}\contentspage}

%% =============================================================================
\begin{document}

% ── Title ─────────────────────────────────────────────────────────────────────
\begin{center}
  {\Huge\bfseries Digital Logic \& Computer Architecture}\\[6pt]
  {\large Hardware, Binary Algebra \& MIPS Assembly Reference}\\[4pt]
  \rule{\textwidth}{0.8pt}
\end{center}

\vspace{8pt}
\tableofcontents
\newpage

%% =============================================================================
%%  PART I – Binary Number Systems & Arithmetic
%% =============================================================================
\greensections
\addcontentsline{toc}{section}{\textcolor{sectiongreen}{PART I \textemdash\ Binary Number Systems \& Arithmetic}}
\addcontentsline{toc}{subsection}{Number Bases Overview}
\addcontentsline{toc}{subsection}{Converting Between Bases}
\addcontentsline{toc}{subsection}{Binary Addition}
\addcontentsline{toc}{subsection}{Two's Complement \& Signed Numbers}
\addcontentsline{toc}{subsection}{Binary Subtraction via Two's Complement}
\addcontentsline{toc}{subsection}{Overflow Detection}
\partbanner{part2bg}{sectiongreen}{Part I \quad Binary Number Systems \& Arithmetic}

\vspace{6pt}
\begin{center}\small\itshape
  Computers only understand 0 and 1. Everything else --- numbers, text, instructions ---
  is encoded in binary. This is where it all starts.
\end{center}
\vspace{4pt}

% ── Number Bases ──────────────────────────────────────────────────────────────
\section*{Number Bases Overview}

A \textbf{positional number system} assigns a value to each digit based on its position (its \emph{weight}). The base determines how many symbols exist and what each position is worth.

\[
  \text{Value} = \sum_{i} d_i \times \text{base}^{\,i}
  \quad \text{(rightmost digit is position } i = 0\text{)}
\]

\noindent
{\renewcommand{\arraystretch}{1.5}
\begin{tabularx}{\textwidth}{>{\bfseries\raggedright\arraybackslash}p{3cm}
                              >{\centering\arraybackslash}p{1.5cm}
                              >{\raggedright\arraybackslash}p{4cm}
                              >{\raggedright\arraybackslash}X}
\toprule
\rowcolor{tableheadergreen}
\textbf{Base} & \textbf{Symbols} & \textbf{Digits Used} & \textbf{Where Used} \\
\midrule
\rowcolor{rowA}  Decimal (10)     & 10 & 0--9 & Human counting; high-level arithmetic \\
\rowcolor{rowBg} Binary (2)       & 2  & 0, 1 & Hardware; every circuit, register, memory cell \\
\rowcolor{rowA}  Octal (8)        & 8  & 0--7 & Older Unix permissions; compact binary shorthand \\
\rowcolor{rowBg} Hexadecimal (16) & 16 & 0--9, A--F & Memory addresses, color codes, machine code \\
\bottomrule
\end{tabularx}}

\vspace{4pt}
\textit{Hex is so common because 1 hex digit = exactly 4 binary bits (a \textbf{nibble}), so 2 hex digits = 1 byte. A 32-bit MIPS address like \texttt{0x00401000} is far more readable than 32 binary digits.}

% ── Converting Between Bases ──────────────────────────────────────────────────
\section*{Converting Between Bases}

\subsection*{Binary $\to$ Decimal}
Multiply each bit by its positional power of 2 and sum:
\[
  1011_2 = 1\times2^3 + 0\times2^2 + 1\times2^1 + 1\times2^0 = 8 + 0 + 2 + 1 = 11_{10}
\]

\subsection*{Decimal $\to$ Binary (Repeated Division by 2)}
Divide by 2 repeatedly; remainders read \textbf{bottom to top} give the binary result:

\[
  43 \div 2 = 21\ R\ \mathbf{1},\quad
  21 \div 2 = 10\ R\ \mathbf{1},\quad
  10 \div 2 = 5\ R\ \mathbf{0},\quad
  5 \div 2 = 2\ R\ \mathbf{1},\quad
  2 \div 2 = 1\ R\ \mathbf{0},\quad
  1 \div 2 = 0\ R\ \mathbf{1}
\]
Reading remainders bottom to top: $43_{10} = 101011_2$

\subsection*{Binary $\leftrightarrow$ Hexadecimal (Group by 4)}
Split binary into groups of 4 from the right; convert each group to one hex digit:
\[
  \underbrace{1010}_{A}\ \underbrace{1111}_{F}\ \underbrace{0011}_{3}{}_{2}
  \quad \Longrightarrow \quad \texttt{0xAF3}
\]

\begin{tcolorbox}[colback=notebg, colframe=noteborder, arc=3pt, boxrule=1pt]
\textbf{Hex digit cheat sheet:}
\begin{center}
\renewcommand{\arraystretch}{1.2}
\begin{tabular}{ccc|ccc|ccc|ccc}
\textbf{Bin} & \textbf{Dec} & \textbf{Hex} &
\textbf{Bin} & \textbf{Dec} & \textbf{Hex} &
\textbf{Bin} & \textbf{Dec} & \textbf{Hex} &
\textbf{Bin} & \textbf{Dec} & \textbf{Hex} \\
\hline
0000 & 0  & 0 & 0100 & 4  & 4 & 1000 & 8  & 8 & 1100 & 12 & C \\
0001 & 1  & 1 & 0101 & 5  & 5 & 1001 & 9  & 9 & 1101 & 13 & D \\
0010 & 2  & 2 & 0110 & 6  & 6 & 1010 & 10 & A & 1110 & 14 & E \\
0011 & 3  & 3 & 0111 & 7  & 7 & 1011 & 11 & B & 1111 & 15 & F \\
\end{tabular}
\end{center}
\end{tcolorbox}

% ── Binary Addition ───────────────────────────────────────────────────────────
\section*{Binary Addition}

Binary addition follows the same column-by-column rules as decimal addition, but with only two symbols. The carry rules are:

\[
  0+0=0,\quad 0+1=1,\quad 1+0=1,\quad 1+1=10_2\ (\text{write 0, carry 1}),\quad 1+1+1=11_2\ (\text{write 1, carry 1})
\]

\begin{tcolorbox}[colback=part2bg, colframe=sectiongreen, arc=3pt, boxrule=1pt]
\textbf{Example:} Add $0110_2\ (6)$ and $0111_2\ (7)$:
\[
\begin{array}{r}
  \textit{carry:}\quad 0\;1\;1\;1\;0 \\
  \phantom{+\ }0\;1\;1\;0 \\
  +\quad 0\;1\;1\;1 \\
  \hline
  \phantom{+\ }1\;1\;0\;1
\end{array}
\]
Result: $1101_2 = 13_{10}$ \checkmark
\end{tcolorbox}

% ── Two's Complement ──────────────────────────────────────────────────────────
\section*{Two's Complement \& Signed Numbers}

Computers represent negative numbers using \textbf{two's complement}. For an $n$-bit number, the \textbf{most significant bit (MSB)} is the \emph{sign bit}: 0 = positive, 1 = negative.

\textbf{Representable range for $n$ bits:}
\[
  -2^{n-1} \quad \text{to} \quad +2^{n-1}-1
  \quad\Longrightarrow\quad
  \text{32-bit:}\quad -2{,}147{,}483{,}648 \text{ to } +2{,}147{,}483{,}647
\]

\textbf{How to negate a number (find two's complement):}
\begin{enumerate}[noitemsep]
  \item \textbf{Invert all bits} (flip every 0 to 1 and every 1 to 0) --- this is the \emph{one's complement}.
  \item \textbf{Add 1} to the result.
\end{enumerate}

\begin{tcolorbox}[colback=part2bg, colframe=sectiongreen, arc=3pt, boxrule=1pt]
\textbf{Example:} Represent $-5$ in 8-bit two's complement.
\begin{enumerate}[noitemsep, leftmargin=*]
  \item Start with $+5$: \quad $00000101$
  \item Invert all bits: \quad $11111010$
  \item Add 1: \quad $11111010 + 1 = \mathbf{11111011}$
\end{enumerate}
So $-5$ is stored as \texttt{11111011}. To verify: $-128 + 64 + 32 + 16 + 8 + 2 + 1 = -128 + 123 = -5$ \checkmark
\end{tcolorbox}

\textbf{Decoding a negative two's complement number:} If the MSB is 1, the value is negative. Apply the same procedure (invert + add 1) to find the magnitude, then negate.

% ── Binary Subtraction ────────────────────────────────────────────────────────
\section*{Binary Subtraction via Two's Complement}

Hardware doesn't need a separate subtractor circuit. Instead, $A - B$ is computed as:
\[
  A - B = A + (-B) = A + (\text{two's complement of } B)
\]

\begin{tcolorbox}[colback=part2bg, colframe=sectiongreen, arc=3pt, boxrule=1pt]
\textbf{Example:} Compute $9 - 5$ in 8-bit two's complement.
\begin{enumerate}[noitemsep, leftmargin=*]
  \item $9 = 00001001$
  \item $-5 = 11111011$ (from earlier)
  \item Add: $00001001 + 11111011 = 1\ 00000100$
  \item Discard the carry out of bit 7: result = $00000100 = 4_{10}$ \checkmark
\end{enumerate}
\end{tcolorbox}

% ── Overflow ──────────────────────────────────────────────────────────────────
\section*{Overflow Detection}

\textbf{Overflow} occurs when the result of an arithmetic operation doesn't fit in the available bits and wraps around to the wrong sign.

\begin{itemize}[noitemsep]
  \item \textbf{Signed overflow rule:} Overflow occurs if and only if the carry \emph{into} the sign bit differs from the carry \emph{out of} the sign bit.
  \item Equivalently: adding two positives gives a negative, or adding two negatives gives a positive.
  \item \textbf{Unsigned overflow:} Simply any carry out of the MSB.
\end{itemize}

\noindent
{\renewcommand{\arraystretch}{1.5}
\begin{tabularx}{\textwidth}{>{\bfseries\raggedright\arraybackslash}p{3.5cm}
                              >{\raggedright\arraybackslash}X
                              >{\centering\arraybackslash}p{2.5cm}}
\toprule
\rowcolor{tableheadergreen}
\textbf{Operation} & \textbf{Overflow occurs when\ldots} & \textbf{Overflow possible?} \\
\midrule
\rowcolor{rowA}  Pos $+$ Pos & Result is negative (MSB = 1) & Yes \\
\rowcolor{rowBg} Neg $+$ Neg & Result is positive (MSB = 0) & Yes \\
\rowcolor{rowA}  Pos $+$ Neg & Never & No \\
\rowcolor{rowBg} Pos $-$ Neg & Same as Pos $+$ Pos & Yes \\
\rowcolor{rowA}  Neg $-$ Pos & Same as Neg $+$ Neg & Yes \\
\bottomrule
\end{tabularx}}

\newpage

%% =============================================================================
%%  PART II – Boolean Algebra & Logic Gates
%% =============================================================================
\orangesections
\addcontentsline{toc}{section}{\textcolor{sectionorange}{PART II \textemdash\ Boolean Algebra \& Logic Gates}}
\addcontentsline{toc}{subsection}{Boolean Algebra Fundamentals}
\addcontentsline{toc}{subsection}{Boolean Algebra Laws \& Identities}
\addcontentsline{toc}{subsection}{The Seven Core Logic Gates}
\addcontentsline{toc}{subsection}{Gate Symbol Standards: ANSI vs.\ IEC}
\addcontentsline{toc}{subsection}{Universal Gates: NAND \& NOR}
\addcontentsline{toc}{subsection}{Boolean Simplification --- Worked Examples}
\addcontentsline{toc}{subsection}{Sum of Products \& Product of Sums}
\partbanner{part4bg}{sectionorange}{Part II \quad Boolean Algebra \& Logic Gates}

\vspace{8pt}

% ── Boolean Algebra Fundamentals ──────────────────────────────────────────────
\section*{Boolean Algebra Fundamentals}

\textbf{Boolean algebra} is the mathematics of logic. Variables can only be \textbf{0} (false) or \textbf{1} (true). Three fundamental operations exist:

\[
  \textbf{AND:}\quad A \cdot B \qquad
  \textbf{OR:}\quad A + B \qquad
  \textbf{NOT:}\quad \overline{A}
\]

Everything in digital logic --- from a single gate to a billion-transistor CPU --- is built from combinations of these three operations.

% ── Boolean Laws ──────────────────────────────────────────────────────────────
\section*{Boolean Algebra Laws \& Identities}

These laws let you simplify and manipulate Boolean expressions, reducing the number of gates needed in a circuit.

\noindent
{\renewcommand{\arraystretch}{1.6}
\begin{tabularx}{\textwidth}{>{\bfseries\raggedright\arraybackslash}p{3.8cm}
                              >{\raggedright\arraybackslash}X
                              >{\raggedright\arraybackslash}X}
\toprule
\rowcolor{tableheaderorange}
\textbf{Law / Identity} & \textbf{AND form} & \textbf{OR form} \\
\midrule
\rowcolor{rowOr} Identity        & $A \cdot 1 = A$              & $A + 0 = A$ \\
\rowcolor{rowA}  Null / Annihilator & $A \cdot 0 = 0$           & $A + 1 = 1$ \\
\rowcolor{rowOr} Idempotent      & $A \cdot A = A$              & $A + A = A$ \\
\rowcolor{rowA}  Complement      & $A \cdot \overline{A} = 0$   & $A + \overline{A} = 1$ \\
\rowcolor{rowOr} Double Negation & \multicolumn{2}{l}{$\overline{\overline{A}} = A$} \\
\rowcolor{rowA}  Commutative     & $A \cdot B = B \cdot A$      & $A + B = B + A$ \\
\rowcolor{rowOr} Associative     & $(A \cdot B) \cdot C = A \cdot (B \cdot C)$ & $(A+B)+C = A+(B+C)$ \\
\rowcolor{rowA}  Distributive    & $A \cdot (B + C) = A\cdot B + A\cdot C$ & $A + (B \cdot C) = (A+B)\cdot(A+C)$ \\
\rowcolor{rowOr} Absorption      & $A \cdot (A + B) = A$        & $A + (A \cdot B) = A$ \\
\rowcolor{rowA}  De Morgan's (1) & \multicolumn{2}{l}{$\overline{A \cdot B} = \overline{A} + \overline{B}$} \\
\rowcolor{rowOr} De Morgan's (2) & \multicolumn{2}{l}{$\overline{A + B} = \overline{A} \cdot \overline{B}$} \\
\bottomrule
\end{tabularx}}

\begin{tcolorbox}[colback=notebg, colframe=noteborder, arc=3pt, boxrule=1pt]
\textbf{De Morgan's Theorem} is the most important law for circuit design. It tells you how to convert between NAND and NOR, and how to ``push a bubble'' through a gate:
\begin{itemize}[noitemsep, leftmargin=*]
  \item \textbf{NAND $\equiv$ OR with inverted inputs:} $\overline{A \cdot B} = \overline{A} + \overline{B}$
  \item \textbf{NOR $\equiv$ AND with inverted inputs:} $\overline{A + B} = \overline{A} \cdot \overline{B}$
\end{itemize}
\textit{Memory: ``Break the bar, change the operation.'' A bar over $A \cdot B$ breaks into $\bar{A} + \bar{B}$. The dot becomes a plus.}
\end{tcolorbox}

% ── The Seven Core Logic Gates ─────────────────────────────────────────────────
\section*{The Seven Core Logic Gates}

\noindent
{\renewcommand{\arraystretch}{1.6}
\begin{tabularx}{\textwidth}{>{\bfseries\raggedright\arraybackslash}p{2cm}
                              >{\centering\arraybackslash}p{2.2cm}
                              >{\raggedright\arraybackslash}p{3.8cm}
                              >{\raggedright\arraybackslash}X}
\toprule
\rowcolor{tableheaderorange}
\textbf{Gate} & \textbf{Expression} & \textbf{Truth Table (AB$\to$Out)} & \textbf{Behavior} \\
\midrule
\rowcolor{rowOr} AND  & $A \cdot B$ & 00$\to$0, 01$\to$0, 10$\to$0, 11$\to$1 & Output is 1 only when \textbf{all} inputs are 1 \\
\rowcolor{rowA}  OR   & $A + B$     & 00$\to$0, 01$\to$1, 10$\to$1, 11$\to$1 & Output is 1 when \textbf{any} input is 1 \\
\rowcolor{rowOr} NOT  & $\overline{A}$ & 0$\to$1, 1$\to$0 & Inverts a single input \\
\rowcolor{rowA}  NAND & $\overline{A \cdot B}$ & 00$\to$1, 01$\to$1, 10$\to$1, 11$\to$0 & AND then invert; output is 0 only when \textbf{all} inputs are 1 \\
\rowcolor{rowOr} NOR  & $\overline{A + B}$     & 00$\to$1, 01$\to$0, 10$\to$0, 11$\to$0 & OR then invert; output is 1 only when \textbf{all} inputs are 0 \\
\rowcolor{rowA}  XOR  & $A \oplus B$           & 00$\to$0, 01$\to$1, 10$\to$1, 11$\to$0 & Output is 1 when inputs are \textbf{different} \\
\rowcolor{rowOr} XNOR & $\overline{A \oplus B}$ & 00$\to$1, 01$\to$0, 10$\to$0, 11$\to$1 & Output is 1 when inputs are \textbf{equal} \\
\bottomrule
\end{tabularx}}

% ── ANSI vs IEC ───────────────────────────────────────────────────────────────
\section*{Gate Symbol Standards: ANSI vs.\ IEC}

\noindent
{\renewcommand{\arraystretch}{1.6}
\begin{tabularx}{\textwidth}{>{\bfseries\raggedright\arraybackslash}p{2cm}
                              >{\raggedright\arraybackslash}X
                              >{\raggedright\arraybackslash}X}
\toprule
\rowcolor{tableheaderorange}
\textbf{Gate} & \textbf{ANSI/IEEE (shapes)} & \textbf{IEC 60617 (boxes)} \\
\midrule
\rowcolor{rowOr} AND  & Flat-back D-shape & Rectangle labeled \textbf{\&} \\
\rowcolor{rowA}  OR   & Curved ``shield'' (pointy output) & Rectangle labeled $\geq\mathbf{1}$ \\
\rowcolor{rowOr} NOT  & Triangle with a bubble at output & Rectangle labeled \textbf{1} with bubble \\
\rowcolor{rowA}  NAND & AND shape + bubble at output & Rectangle labeled \textbf{\&} + bubble \\
\rowcolor{rowOr} NOR  & OR shape + bubble at output & Rectangle labeled $\geq\mathbf{1}$ + bubble \\
\rowcolor{rowA}  XOR  & OR shape with extra curved input line & Rectangle labeled $=\mathbf{1}$ \\
\rowcolor{rowOr} XNOR & XOR shape + bubble at output & Rectangle labeled $=\mathbf{1}$ + bubble \\
\bottomrule
\end{tabularx}}

\vspace{4pt}
\begin{tcolorbox}[colback=part4bg, colframe=sectionorange, arc=3pt, boxrule=1pt]
\textbf{The Bubble Convention:} A small circle (\emph{bubble}) placed at any input or output indicates \textbf{inversion (NOT)}. This is how NAND, NOR, and XNOR are drawn --- take the base gate and add a bubble at the output. A bubble at an \emph{input} means that input is active-low (the signal is logically true when the wire is at 0V).
\end{tcolorbox}

% ── Universal Gates ───────────────────────────────────────────────────────────
\section*{Universal Gates: NAND \& NOR}

\textbf{NAND and NOR are each individually ``universal''} --- any Boolean function can be built using \emph{only} NAND gates (or only NOR gates). This matters for chip manufacturing because a fab only needs one gate type.

\textbf{Building AND, OR, NOT from NAND only:}
\[
  \textbf{NOT:}\quad \overline{A} = \overline{A \cdot A} \quad\text{(NAND with both inputs tied together)}
\]
\[
  \textbf{AND:}\quad A \cdot B = \overline{\overline{A \cdot B}} \quad\text{(NAND then NOT = NAND-NAND)}
\]
\[
  \textbf{OR:}\quad A + B = \overline{\overline{A} \cdot \overline{B}} \quad\text{(De Morgan: NOT each input, then NAND)}
\]

% ── Boolean Simplification ────────────────────────────────────────────────────
\section*{Boolean Simplification --- Worked Examples}

\begin{tcolorbox}[colback=part4bg, colframe=sectionorange,
  title={\bfseries Example 1: Simplify $F = A\cdot B + A\cdot\overline{B}$}, arc=3pt, boxrule=1pt]
\begin{align*}
  F &= A\cdot B + A\cdot\overline{B} \\
    &= A\cdot(B + \overline{B}) \quad\text{(Distributive)} \\
    &= A \cdot 1               \quad\text{(Complement law: $B + \overline{B} = 1$)} \\
    &= A                        \quad\text{(Identity)}
\end{align*}
\textit{The expression simplifies to just $A$. The $B$ input is irrelevant --- this circuit can be replaced by a direct wire from $A$.}
\end{tcolorbox}

\begin{tcolorbox}[colback=part4bg, colframe=sectionorange,
  title={\bfseries Example 2: Simplify $F = \overline{A}\cdot B + A\cdot B + A\cdot\overline{B}$}, arc=3pt, boxrule=1pt]
\begin{align*}
  F &= \overline{A}\cdot B + A\cdot B + A\cdot\overline{B} \\
    &= B(\overline{A} + A) + A\cdot\overline{B}            \quad\text{(Distributive)} \\
    &= B\cdot 1 + A\cdot\overline{B}                       \quad\text{(Complement)} \\
    &= B + A\cdot\overline{B}                              \quad\text{(Identity)} \\
    &= (B + A)\cdot(B + \overline{B})                      \quad\text{(Distributive OR form)} \\
    &= (B + A)\cdot 1 = A + B
\end{align*}
\textit{Three AND gates and an OR reduced to a single OR gate.}
\end{tcolorbox}

% ── SOP and POS ───────────────────────────────────────────────────────────────
\section*{Sum of Products (SOP) \& Product of Sums (POS)}

These are the two canonical forms for expressing any Boolean function from a truth table.

\begin{itemize}[noitemsep]
  \item \textbf{Sum of Products (SOP):} Take every row where the output is \textbf{1}. For each such row, AND the inputs together (a \emph{minterm}), then OR all the minterms.
  \item \textbf{Product of Sums (POS):} Take every row where the output is \textbf{0}. For each such row, OR the inputs together (a \emph{maxterm}), then AND all the maxterms.
\end{itemize}

\begin{tcolorbox}[colback=part4bg, colframe=sectionorange,
  title={\bfseries Example: Build an SOP expression for XOR from a truth table}, arc=3pt, boxrule=1pt]
\begin{center}
\renewcommand{\arraystretch}{1.3}
\begin{tabular}{cc|c|l}
$A$ & $B$ & $A \oplus B$ & Minterm \\
\hline
0 & 0 & 0 & (skip) \\
0 & 1 & 1 & $\overline{A}\cdot B$ \\
1 & 0 & 1 & $A\cdot\overline{B}$ \\
1 & 1 & 0 & (skip) \\
\end{tabular}
\end{center}
\[
  F = \overline{A}\cdot B + A\cdot\overline{B}
\]
This is exactly the textbook definition of XOR --- now you know how to derive it from scratch.
\end{tcolorbox}

\newpage

%% =============================================================================
%%  PART III – Combinational Circuits
%% =============================================================================
\tealsections
\addcontentsline{toc}{section}{\textcolor{sectionteal}{PART III \textemdash\ Combinational Circuits}}
\addcontentsline{toc}{subsection}{Half Adder}
\addcontentsline{toc}{subsection}{Full Adder}
\addcontentsline{toc}{subsection}{Ripple Carry Adder}
\addcontentsline{toc}{subsection}{Multiplexer (MUX)}
\addcontentsline{toc}{subsection}{Decoder}
\addcontentsline{toc}{subsection}{The Arithmetic Logic Unit (ALU)}
\partbanner{part5bg}{sectionteal}{Part III \quad Combinational Circuits}

\vspace{8pt}
\begin{center}\small\itshape
  Combinational circuits produce outputs that depend \emph{only} on current inputs --- no memory, no feedback.
\end{center}

% ── Half Adder ────────────────────────────────────────────────────────────────
\section*{Half Adder}

A \textbf{half adder} adds two single bits and produces a \textbf{Sum} and a \textbf{Carry}:

\[
  \text{Sum} = A \oplus B \qquad \text{Carry} = A \cdot B
\]

\begin{center}
\renewcommand{\arraystretch}{1.3}
\begin{tabular}{cc|cc}
\toprule
$A$ & $B$ & Sum & Carry \\
\midrule
0 & 0 & 0 & 0 \\
0 & 1 & 1 & 0 \\
1 & 0 & 1 & 0 \\
1 & 1 & 0 & 1 \\
\bottomrule
\end{tabular}
\end{center}

It's called ``half'' because it can't accept a carry-in from a previous column --- that's what the full adder adds.

% ── Full Adder ────────────────────────────────────────────────────────────────
\section*{Full Adder}

A \textbf{full adder} adds three bits: $A$, $B$, and $C_{\mathrm{in}}$ (carry-in from previous column):

\[
  \text{Sum} = A \oplus B \oplus C_{\mathrm{in}}
  \qquad
  C_{\mathrm{out}} = (A \cdot B) + (C_{\mathrm{in}} \cdot (A \oplus B))
\]

A full adder can be built from two half adders plus an OR gate:

\begin{center}
\renewcommand{\arraystretch}{1.3}
\begin{tabular}{ccc|cc}
\toprule
$A$ & $B$ & $C_{\mathrm{in}}$ & Sum & $C_{\mathrm{out}}$ \\
\midrule
0&0&0 & 0 & 0 \\
0&0&1 & 1 & 0 \\
0&1&0 & 1 & 0 \\
0&1&1 & 0 & 1 \\
1&0&0 & 1 & 0 \\
1&0&1 & 0 & 1 \\
1&1&0 & 0 & 1 \\
1&1&1 & 1 & 1 \\
\bottomrule
\end{tabular}
\end{center}

% ── Ripple Carry Adder ────────────────────────────────────────────────────────
\section*{Ripple Carry Adder}

Chain $n$ full adders together and you get an \textbf{$n$-bit ripple carry adder}. The $C_{\mathrm{out}}$ of each stage feeds into the $C_{\mathrm{in}}$ of the next. The first stage's $C_{\mathrm{in}}$ is tied to 0.

\[
  A_3A_2A_1A_0 + B_3B_2B_1B_0
  \quad\xrightarrow{\text{4 full adders}}\quad
  S_3S_2S_1S_0,\quad C_{\mathrm{out}}
\]

\begin{tcolorbox}[colback=part5bg, colframe=sectionteal, arc=3pt, boxrule=1pt]
\textbf{Limitation --- the carry ripple delay:} The final sum bits can't be computed until the carry has propagated through every stage from right to left. For a 32-bit adder, this means 32 gate delays. This is why modern CPUs use \textbf{carry-lookahead adders} that compute carries in parallel, reducing delay to $O(\log n)$.
\end{tcolorbox}

% ── Multiplexer ───────────────────────────────────────────────────────────────
\section*{Multiplexer (MUX)}

A \textbf{multiplexer} selects one of $2^n$ inputs based on $n$ \textbf{select} lines and routes it to the single output. Think of it as a programmable switch.

\[
  \text{2-to-1 MUX:} \quad Y = S\cdot B + \overline{S}\cdot A
\]

When $S=0$, output is $A$; when $S=1$, output is $B$. MUXes are used everywhere in a CPU: choosing between two ALU operands, selecting between PC+4 and a branch target, routing data through datapaths.

% ── Decoder ───────────────────────────────────────────────────────────────────
\section*{Decoder}

An $n$-to-$2^n$ \textbf{decoder} takes $n$ input bits and asserts exactly \textbf{one} of its $2^n$ output lines high, corresponding to the binary value of the input. Used to select memory locations, registers, or instruction types.

\textit{Example: a 2-to-4 decoder with input $AB = 10$ asserts output line 2 (index from 0).}

% ── ALU ───────────────────────────────────────────────────────────────────────
\section*{The Arithmetic Logic Unit (ALU)}

The \textbf{ALU} is the computational core of the CPU. It takes two operands and a \textbf{control signal} that selects which operation to perform:

\noindent
{\renewcommand{\arraystretch}{1.5}
\begin{tabularx}{\textwidth}{>{\centering\arraybackslash}p{2.5cm}
                              >{\bfseries\centering\arraybackslash}p{2.5cm}
                              >{\raggedright\arraybackslash}X}
\toprule
\rowcolor{tableheaderteal}
\textbf{ALU Control} & \textbf{Operation} & \textbf{Used by MIPS instruction} \\
\midrule
\rowcolor{rowTe} 0000 & AND & \texttt{and} \\
\rowcolor{rowA}  0001 & OR  & \texttt{or} \\
\rowcolor{rowTe} 0010 & Add & \texttt{add}, \texttt{lw}, \texttt{sw}, \texttt{beq} \\
\rowcolor{rowA}  0110 & Subtract & \texttt{sub}, \texttt{beq} (subtract and test zero) \\
\rowcolor{rowTe} 0111 & Set-on-Less-Than & \texttt{slt} \\
\rowcolor{rowA}  1100 & NOR & \texttt{nor} \\
\bottomrule
\end{tabularx}}

\newpage

%% =============================================================================
%%  PART IV – MIPS Assembly Language
%% =============================================================================
\bluesections
\addcontentsline{toc}{section}{\textcolor{sectionblue}{PART IV \textemdash\ MIPS Assembly Language}}
\addcontentsline{toc}{subsection}{Why Assembly?}
\addcontentsline{toc}{subsection}{MIPS Register File}
\addcontentsline{toc}{subsection}{Instruction Format Overview}
\addcontentsline{toc}{subsection}{R-Type Instructions (Register-to-Register)}
\addcontentsline{toc}{subsection}{I-Type Instructions (Immediate \& Memory)}
\addcontentsline{toc}{subsection}{J-Type Instructions (Jump)}
\addcontentsline{toc}{subsection}{Memory Addressing: Md[address]}
\addcontentsline{toc}{subsection}{Control Flow: Branches \& Jumps}
\addcontentsline{toc}{subsection}{Functions \& the Call Convention}
\partbanner{partbg}{sectionblue}{Part IV \quad MIPS Assembly Language}

\vspace{8pt}

% ── Why Assembly ──────────────────────────────────────────────────────────────
\section*{Why Assembly?}

Assembly is the \textbf{lowest human-readable level} of programming. Each instruction maps directly to a single machine instruction that the CPU executes. Understanding it means understanding exactly what the hardware does.

\begin{itemize}[noitemsep]
  \item \textbf{MIPS} (Microprocessor without Interlocked Pipeline Stages) is a \textbf{RISC} (Reduced Instruction Set Computer) architecture. It has a small, regular set of simple instructions.
  \item All MIPS instructions are exactly \textbf{32 bits} wide --- no variable-length encoding.
  \item \textbf{Load/store architecture:} Only \texttt{lw} and \texttt{sw} touch memory. All arithmetic operates on \emph{registers} only.
\end{itemize}

% ── Register File ─────────────────────────────────────────────────────────────
\section*{MIPS Register File}

MIPS has \textbf{32 registers}, each 32 bits wide. They are addressed by a 5-bit number (0--31) in hardware, but referred to by name in assembly. Think of the register file as a tiny, ultra-fast array of 32 slots that lives inside the CPU.

\noindent
{\small\renewcommand{\arraystretch}{1.5}
\begin{tabularx}{\textwidth}{>{\ttfamily\centering\arraybackslash}p{2.4cm}
                              >{\centering\arraybackslash}p{1.4cm}
                              >{\raggedright\arraybackslash}X
                              >{\centering\arraybackslash}p{2.2cm}}
\toprule
\rowcolor{tableheader}
\textbf{Name} & \textbf{Number} & \textbf{Purpose} & \textbf{Preserved?} \\
\midrule
\rowcolor{rowA}  \$zero & 0      & Hardwired to 0; writes are ignored & N/A \\
\rowcolor{rowB}  \$at   & 1      & Assembler temporary; do not use directly & No \\
\rowcolor{rowA}  \$v0--\$v1 & 2--3 & Function return values & No \\
\rowcolor{rowB}  \$a0--\$a3 & 4--7 & Function arguments (first 4) & No \\
\rowcolor{rowA}  \$t0--\$t7 & 8--15 & Temporaries; caller-saved (may be clobbered by a called function) & \textbf{No} \\
\rowcolor{rowB}  \$s0--\$s7 & 16--23 & Saved registers; callee-saved (a called function must restore these) & \textbf{Yes} \\
\rowcolor{rowA}  \$t8--\$t9 & 24--25 & Additional temporaries & No \\
\rowcolor{rowB}  \$gp & 28 & Global pointer & Yes \\
\rowcolor{rowA}  \$sp & 29 & Stack pointer; points to top of current stack frame & Yes \\
\rowcolor{rowB}  \$fp & 30 & Frame pointer; optional base of current stack frame & Yes \\
\rowcolor{rowA}  \$ra & 31 & Return address; set by \texttt{jal} & Yes \\
\bottomrule
\end{tabularx}}

\begin{tcolorbox}[colback=partbg, colframe=sectionblue, arc=3pt, boxrule=1pt]
\textbf{\$zero's many uses:} Since \$zero is always 0, it is used as a constant operand to fake instructions: \\
\texttt{move \$t0, \$t1} $\equiv$ \texttt{add \$t0, \$t1, \$zero} \quad (add zero to copy a register)\\
\texttt{clear \$t0} $\equiv$ \texttt{add \$t0, \$zero, \$zero} \quad (add two zeros)\\
\texttt{beq \$t0, \$zero, label} \quad (branch if \$t0 equals zero --- test for zero)
\end{tcolorbox}

% ── Instruction Format Overview ───────────────────────────────────────────────
\section*{Instruction Format Overview}

Every MIPS instruction is encoded as a \textbf{32-bit binary word}. There are three formats:

\[
  \underbrace{\textbf{R-type}}_{\text{Register ops}}
  \qquad
  \underbrace{\textbf{I-type}}_{\text{Immediate / Memory}}
  \qquad
  \underbrace{\textbf{J-type}}_{\text{Jumps}}
\]

\noindent
{\footnotesize\renewcommand{\arraystretch}{1.5}
\begin{tabularx}{\textwidth}{>{\bfseries\centering\arraybackslash}p{1.2cm}
                              >{\centering\arraybackslash}p{1.5cm}
                              >{\centering\arraybackslash}p{1.5cm}
                              >{\centering\arraybackslash}p{1.5cm}
                              >{\centering\arraybackslash}p{1.5cm}
                              >{\centering\arraybackslash}p{1.5cm}
                              >{\raggedright\arraybackslash}X}
\toprule
\rowcolor{tableheader}
\textbf{Type} & \textbf{op [31:26]} & \textbf{rs [25:21]} & \textbf{rt [20:16]} & \textbf{rd [15:11]} & \textbf{shamt [10:6]} & \textbf{funct / imm / target [5:0 or more]} \\
\midrule
\rowcolor{rowA} R & 6 bits (000000) & 5 bits & 5 bits & 5 bits & 5 bits & funct: 6 bits \\
\rowcolor{rowB} I & 6 bits & 5 bits (rs) & 5 bits (rt) & \multicolumn{3}{l}{immediate: 16 bits} \\
\rowcolor{rowA} J & 6 bits & \multicolumn{5}{l}{target address: 26 bits} \\
\bottomrule
\end{tabularx}}

% ── R-Type ────────────────────────────────────────────────────────────────────
\section*{R-Type Instructions (Register-to-Register)}

\textbf{R-type} is the most common format. It performs an operation on two source registers and writes the result to a destination register.

\begin{tcolorbox}[colback=partbg, colframe=sectionblue, arc=3pt, boxrule=1pt]
\textbf{Syntax pattern:} \quad \texttt{op\ \ rd,\ rs,\ rt}\\[4pt]
\texttt{rd} = \textbf{destination} (where the result goes)\\
\texttt{rs} = \textbf{source 1} (first operand)\\
\texttt{rt} = \textbf{source 2} (second operand)\\[4pt]
Think of it as: \textbf{(result, value1, value2)} --- the destination always comes first.
\end{tcolorbox}

\begin{lstlisting}[style=mips, caption={R-Type examples}]
add   $t0, $s1, $s2   # $t0 = $s1 + $s2
sub   $t1, $s3, $s4   # $t1 = $s3 - $s4
and   $t2, $s0, $s1   # $t2 = $s0 AND $s1  (bitwise)
or    $t3, $s0, $s1   # $t3 = $s0 OR  $s1  (bitwise)
nor   $t4, $s0, $zero # $t4 = NOT $s0  (NOR with zero = NOT)
slt   $t5, $s1, $s2   # $t5 = 1 if $s1 < $s2, else 0
sll   $t6, $s0, 2     # $t6 = $s0 << 2  (shift left by 2 bits = multiply by 4)
srl   $t7, $s0, 2     # $t7 = $s0 >> 2  (shift right = divide by 4)
\end{lstlisting}

% ── I-Type ────────────────────────────────────────────────────────────────────
\section*{I-Type Instructions (Immediate \& Memory)}

\textbf{I-type} encodes a 16-bit constant (\emph{immediate}) directly in the instruction --- no need to load it into a register first. It is also used for all memory access.

\begin{tcolorbox}[colback=partbg, colframe=sectionblue, arc=3pt, boxrule=1pt]
\textbf{Arithmetic immediate syntax:} \quad \texttt{op\ \ rt,\ rs,\ imm}\\[2pt]
\texttt{rt} = destination; \quad \texttt{rs} = source register; \quad \texttt{imm} = constant (signed 16-bit)\\[6pt]
\textbf{Memory (load/store) syntax:} \quad \texttt{op\ \ rt,\ offset(rs)}\\[2pt]
\texttt{rs} = \textbf{base register} (holds a memory address); \quad \texttt{offset} = a signed 16-bit constant added to the base\\
The effective memory address is: \textbf{Md[rs + offset]}
\end{tcolorbox}

\begin{lstlisting}[style=mips, caption={I-Type examples}]
addi  $t0, $s1, 10    # $t0 = $s1 + 10  (add immediate)
andi  $t1, $s2, 0xFF  # $t1 = $s2 AND 0xFF  (mask lower 8 bits)
ori   $t2, $zero, 42  # $t2 = 42  (load small constant)
slti  $t3, $s1, 100   # $t3 = 1 if $s1 < 100, else 0

# Memory: lw/sw -- the offset is added to the base register address
lw    $t0, 0($s1)     # $t0 = Md[$s1 + 0]  -- load word at address in $s1
lw    $t1, 4($s1)     # $t1 = Md[$s1 + 4]  -- next word (4 bytes later)
lw    $t2, 8($s1)     # $t2 = Md[$s1 + 8]  -- two words ahead
sw    $t0, 12($s1)    # Md[$s1 + 12] = $t0  -- store word
\end{lstlisting}

% ── Memory Addressing ─────────────────────────────────────────────────────────
\section*{Memory Addressing: Md[address]}

MIPS memory is a \textbf{giant array of bytes}, each with a unique 32-bit address. You can think of it exactly like an array in a high-level language:

\[
  \texttt{Md[address]} \quad \longleftrightarrow \quad \texttt{array[index]}
\]

\begin{tcolorbox}[colback=partbg, colframe=sectionblue,
  title={\bfseries The Memory Array Analogy}, arc=3pt, boxrule=1pt]
Imagine memory as: \quad \texttt{int Md[4294967296]} \quad (a $2^{32}$-element array of bytes)\\[4pt]
\texttt{Md[0x00401000]} means: ``go to address \texttt{0x00401000} and read the byte there.''\\[4pt]
\texttt{lw} reads \textbf{4 consecutive bytes} (one word) starting at the given address.\\
\texttt{sw} writes \textbf{4 bytes} to that address.\\[4pt]
\textbf{Important:} Words must be \textbf{word-aligned} --- the address must be a multiple of 4. \texttt{Md[5564]} is valid ($5564 = 4 \times 1391$); \texttt{Md[5565]} would cause an alignment fault.
\end{tcolorbox}

\textbf{The offset model} lets you navigate structures and arrays:

\begin{lstlisting}[style=mips, caption={Accessing an array: int A[5] where \$s0 holds the base address of A}]
# A[0] is at Md[$s0 + 0]
# A[1] is at Md[$s0 + 4]   (each int = 4 bytes, so index * 4)
# A[2] is at Md[$s0 + 8]
# A[i] is at Md[$s0 + i*4]

lw    $t0, 0($s0)     # $t0 = A[0]
lw    $t1, 4($s0)     # $t1 = A[1]
lw    $t2, 8($s0)     # $t2 = A[2]
sw    $t3, 12($s0)    # A[3] = $t3

# To access A[i] where i is in $t4:
sll   $t5, $t4, 2     # $t5 = i * 4  (shift left 2 = multiply by 4)
add   $t6, $s0, $t5   # $t6 = base + offset = address of A[i]
lw    $t7, 0($t6)     # $t7 = A[i]
\end{lstlisting}

% ── J-Type ────────────────────────────────────────────────────────────────────
\section*{J-Type Instructions (Jump)}

\textbf{J-type} encodes a 26-bit target address for unconditional jumps. The full 32-bit PC is reconstructed by combining the upper 4 bits of PC with the 26-bit field shifted left by 2.

\begin{lstlisting}[style=mips, caption={J-Type examples}]
j     label            # Unconditional jump to label
jal   function         # Jump and link: jump to function, save PC+4 in $ra
jr    $ra              # Jump register: return from function (jump to address in $ra)
\end{lstlisting}

% ── Control Flow ──────────────────────────────────────────────────────────────
\section*{Control Flow: Branches \& Jumps}

Branches are \textbf{I-type}. The 16-bit immediate is a \textbf{PC-relative offset} in words (signed), added to PC+4 when the branch is taken.

\begin{lstlisting}[style=mips, caption={Branch examples and C-to-MIPS translation}]
# C: if (a == b) goto L
beq   $s0, $s1, L     # branch if $s0 == $s1

# C: if (a != b) goto L
bne   $s0, $s1, L     # branch if $s0 != $s1

# C: if (a < b) goto L
slt   $t0, $s0, $s1   # $t0 = 1 if $s0 < $s1
bne   $t0, $zero, L   # branch if $t0 != 0

# C: for (i = 0; i < 10; i++) { ... }
      addi  $t0, $zero, 0   # i = 0
loop: slti  $t1, $t0, 10    # $t1 = (i < 10)
      beq   $t1, $zero, end # if !(i < 10), exit loop
      # ... loop body ...
      addi  $t0, $t0, 1     # i++
      j     loop
end:
\end{lstlisting}

% ── Functions & Call Convention ───────────────────────────────────────────────
\section*{Functions \& the Call Convention}

\begin{tcolorbox}[colback=partbg, colframe=sectionblue,
  title={\bfseries The Standard MIPS Call Convention}, arc=3pt, boxrule=1pt]
\textbf{Before the call (caller's job):}
\begin{enumerate}[noitemsep, leftmargin=*]
  \item Put arguments in \texttt{\$a0--\$a3} (up to 4 args). Extra args go on the stack.
  \item Execute \texttt{jal function} --- this jumps to the function \emph{and} saves PC+4 into \texttt{\$ra}.
\end{enumerate}
\textbf{Inside the function (callee's job):}
\begin{enumerate}[noitemsep, leftmargin=*]
  \item Save any \texttt{\$s} registers you will use (push them onto the stack).
  \item Save \texttt{\$ra} if you will call another function (otherwise \texttt{jal} would overwrite it).
  \item Do the work; put return value(s) in \texttt{\$v0}/\texttt{\$v1}.
  \item Restore all saved registers from the stack.
  \item Execute \texttt{jr \$ra} to return to the caller.
\end{enumerate}
\end{tcolorbox}

\begin{lstlisting}[style=mips, caption={Complete function example: int square(int n) \{ return n*n; \}}]
# Caller:
addi  $a0, $zero, 7   # argument: n = 7
jal   square           # call square; $ra = address of next instruction
# $v0 now holds 49

# Function:
square:
  mult  $a0, $a0        # HI:LO = n * n  (32x32 -> 64-bit result)
  mflo  $v0             # move lower 32 bits of product to $v0
  jr    $ra             # return to caller
\end{lstlisting}

\begin{lstlisting}[style=mips, caption={Stack frame example for a function that calls another function}]
factorial:
  addi  $sp, $sp, -8    # Allocate 2 words on stack (grow downward)
  sw    $ra, 4($sp)     # Save return address
  sw    $a0, 0($sp)     # Save argument (n)

  slti  $t0, $a0, 1     # $t0 = (n < 1)
  bne   $t0, $zero, base_case

  addi  $a0, $a0, -1    # n - 1
  jal   factorial        # recursive call: factorial(n-1)
  lw    $a0, 0($sp)     # restore n
  mult  $v0, $a0         # n * factorial(n-1)
  mflo  $v0

  lw    $ra, 4($sp)     # restore return address
  addi  $sp, $sp, 8     # deallocate stack frame
  jr    $ra

base_case:
  addi  $v0, $zero, 1   # return 1
  lw    $ra, 4($sp)
  addi  $sp, $sp, 8
  jr    $ra
\end{lstlisting}

\newpage

%% =============================================================================
%%  PART IVb -- MIPS Instruction Reference Tables
%% =============================================================================
\bluesections
\addcontentsline{toc}{section}{\textcolor{sectionblue}{PART IVb \textemdash\ MIPS Instruction Reference Tables}}
\addcontentsline{toc}{subsection}{Complete Instruction Summary}
\addcontentsline{toc}{subsection}{Machine Encoding Reference}
\partbanner{partbg}{sectionblue}{Part IVb \quad MIPS Instruction Reference Tables}

\vspace{6pt}
\begin{center}\small\itshape
  Complete encoding reference for MIPSzy instructions.
  Use alongside the assembly examples in Part IV.
\end{center}

% -- Complete Instruction Summary ---------------------------------------------
\section*{Complete Instruction Summary}

\noindent
{\small\renewcommand{\arraystretch}{1.6}
\begin{tabularx}{\textwidth}{>{\ttfamily\bfseries\raggedright\arraybackslash}p{1.4cm}
                              >{\ttfamily\raggedright\arraybackslash}p{3.6cm}
                              >{\raggedright\arraybackslash}X
                              >{\ttfamily\raggedright\arraybackslash}p{3.2cm}}
\toprule
\rowcolor{tableheader}
\textbf{Instr} & \textbf{Format} & \textbf{Description} & \textbf{Example} \\
\midrule
\rowcolor{rowA}  lw   & lw \$a, 0(\$b)      & Load word: copy from memory at address \texttt{\$b} into register \texttt{\$a}.                                                               & lw \$t3, 0(\$t6)         \\
\rowcolor{rowB}  lw   & lw \$a, C(\$b)      & Load word with offset: copy from address \texttt{\$b + C} into \texttt{\$a}.                                                              & lw \$t3, 20(\$t6)        \\
\rowcolor{rowA}  sw   & sw \$a, 0(\$b)      & Store word: copy register \texttt{\$a} to memory at address \texttt{\$b}.                                                                  & sw \$t1, 0(\$t3)         \\
\rowcolor{rowB}  sw   & sw \$a, C(\$b)      & Store word with offset: copy \texttt{\$a} to address \texttt{\$b + C}.                                                                    & sw \$t1, -4(\$t3)        \\
\rowcolor{rowA}  addi & addi \$a, \$b, C    & Add immediate: \texttt{\$a = \$b + C}. Constant C is sign-extended.                                                                        & addi \$t3, \$t2, 7       \\
\rowcolor{rowB}  add  & add \$a, \$b, \$c  & Add registers: \texttt{\$a = \$b + \$c}.                                                                                                  & add \$t4, \$t1, \$t2     \\
\rowcolor{rowA}  sub  & sub \$a, \$b, \$c  & Subtract: \texttt{\$a = \$b - \$c}.                                                                                                       & sub \$t3, \$t2, \$t5     \\
\rowcolor{rowB}  mul  & mul \$a, \$b, \$c  & Multiply (pseudo): \texttt{\$a} = lower 32 bits of \texttt{\$b * \$c}. Implemented via \texttt{mult} + \texttt{mflo} by the assembler.      & mul \$t3, \$t2, \$t1     \\
\rowcolor{rowA}  mult & mult \$a, \$b      & Multiply (hardware): 64-bit result stored in special registers \texttt{HI:\$LO}.                                                           & mult \$t3, \$t5          \\
\rowcolor{rowB}  mflo & mflo \$a           & Move from LO: copy the lower 32 bits of a \texttt{mult}/\texttt{div} result into \texttt{\$a}.                                            & mflo \$t2                \\
\rowcolor{rowA}  beq  & beq \$a, \$b, Lbl  & Branch if equal: jump to \texttt{Lbl} if \texttt{\$a == \$b}; otherwise fall through.                                                     & beq \$t3, \$t2, SumEq5   \\
\rowcolor{rowB}  bne  & bne \$a, \$b, Lbl  & Branch if not equal: jump to \texttt{Lbl} if \texttt{\$a != \$b}.                                                                         & bne \$t4, \$t5, NeqCorr  \\
\rowcolor{rowA}  slt  & slt \$a, \$b, \$c  & Set on less than: \texttt{\$a = 1} if \texttt{\$b < \$c}, else \texttt{\$a = 0}.                                                          & slt \$t1, \$t5, \$t6     \\
\rowcolor{rowB}  j    & j JLabel           & Jump unconditionally to \texttt{JLabel}.                                                                                                   & j CalcTip                \\
\rowcolor{rowA}  jal  & jal JLabel         & Jump and link: save \texttt{PC+4} into \texttt{\$ra}, then jump to \texttt{JLabel}.                                                       & jal CalcTip              \\
\rowcolor{rowB}  jr   & jr \$a             & Jump register: set \texttt{PC = \$a}. Used for function returns (\texttt{jr \$ra}).                                                       & jr \$ra                  \\
\bottomrule
\end{tabularx}}

\vspace{6pt}

\begin{tcolorbox}[colback=partbg, colframe=sectionblue, arc=3pt, boxrule=1pt]
\textbf{Quick encoding reminder:}
\begin{itemize}[noitemsep, leftmargin=*]
  \item \textbf{R-type} (\texttt{add}, \texttt{sub}, \texttt{slt}, \texttt{mult}, \texttt{mflo}, \texttt{jr}): opcode is always \texttt{000000}; operation selected by the \textbf{funct} field.
  \item \textbf{I-type} (\texttt{lw}, \texttt{sw}, \texttt{addi}, \texttt{beq}, \texttt{bne}): encodes one register source, one destination/target register, and a 16-bit immediate.
  \item \textbf{J-type} (\texttt{j}, \texttt{jal}): encodes a 26-bit target address; the assembler inserts the correct bit pattern.
  \item \textbf{mul} is a \emph{pseudo-instruction} -- the assembler expands it into \texttt{mult} + \texttt{mflo} automatically.
\end{itemize}
\end{tcolorbox}

% -- Machine Encoding Reference ----------------------------------------------
\section*{Machine Encoding Reference}

Each MIPS instruction is exactly 32 bits. The bit-field layout differs by type (see Part IV for the full field breakdown). The table below shows the binary encoding of common instructions.

\noindent
{\small\renewcommand{\arraystretch}{1.6}
\begin{tabularx}{\textwidth}{>{\ttfamily\raggedright\arraybackslash}p{5cm}
                              >{\ttfamily\raggedright\arraybackslash}X
                              >{\centering\arraybackslash}p{1.4cm}}
\toprule
\rowcolor{tableheader}
\textbf{Assembly} & \textbf{Machine Encoding (binary)} & \textbf{Type} \\
\midrule
\rowcolor{rowA}  lw \$t0, 0(\$t1)      & 100011\ \ 01001\ 01000\ \ \ \ 0000000000000000 & I \\
\rowcolor{rowB}  sw \$t0, 0(\$t1)      & 101011\ \ 01001\ 01000\ \ \ \ 0000000000000000 & I \\
\rowcolor{rowA}  addi \$t0, \$t1, 15   & 001000\ \ 01001\ 01000\ \ \ \ 0000000000001111 & I \\
\rowcolor{rowB}  add \$t0, \$t1, \$t2  & 000000\ \ 01001\ 01010\ 01000\ 00000\ 100000  & R \\
\rowcolor{rowA}  sub \$t0, \$t1, \$t2  & 000000\ \ 01001\ 01010\ 01000\ 00000\ 100010  & R \\
\rowcolor{rowB}  mult \$t1, \$t2       & 000000\ \ 01001\ 01010\ 00000\ 00000\ 011000  & R \\
\rowcolor{rowA}  mflo \$t0             & 000000\ \ 00000\ 00000\ 01000\ 00000\ 010010  & R \\
\rowcolor{rowB}  beq \$t1, \$t2, Lbl   & 000100\ \ 01001\ 01010\ \ \ \ 0000000000000010 & I \\
\rowcolor{rowA}  bne \$t1, \$t2, Lbl   & 000101\ \ 01001\ 01010\ \ \ \ 0000000000000010 & I \\
\rowcolor{rowB}  slt \$t0, \$t1, \$t2  & 000000\ \ 01001\ 01010\ 01000\ 00000\ 101010  & R \\
\rowcolor{rowA}  j JLabel              & 000010\ \ \ \ 00000000000000000000000101       & J \\
\rowcolor{rowB}  jal JLabel            & 000011\ \ \ \ 00000000000000100000000101       & J \\
\rowcolor{rowA}  jr \$t1               & 000000\ \ 01001\ 00000\ 00000\ 00000\ 001000  & R \\
\bottomrule
\end{tabularx}}

\vspace{4pt}
\begin{tcolorbox}[colback=notebg, colframe=noteborder, arc=3pt, boxrule=1pt]
\textbf{Reading the encoding:} Fields are separated by spaces for clarity.
For \textbf{R-type}: \texttt{[op:6][rs:5][rt:5][rd:5][shamt:5][funct:6]}.
For \textbf{I-type}: \texttt{[op:6][rs:5][rt:5][immediate:16]}.
For \textbf{J-type}: \texttt{[op:6][target:26]}.
The opcode for all R-type instructions is \texttt{000000}; the \textbf{funct} field distinguishes them.
\end{tcolorbox}

\newpage

%% =============================================================================
%%  PART V -- Quick Reference Cheat Sheet
%% =============================================================================
\purplesections
\addcontentsline{toc}{section}{\textcolor{sectionpurple}{PART V \textemdash\ Quick Reference Cheat Sheet}}
\partbanner{part3bg}{sectionpurple}{Part V \quad Quick Reference Cheat Sheet}

\vspace{4pt}
\begin{center}\small\itshape One page. Everything essential. Great for last-minute review.\end{center}
\vspace{4pt}

% ── Binary & Two's Complement ─────────────────────────────────────────────────
\begin{tcolorbox}[colback=part2bg, colframe=sectiongreen,
  title={\bfseries Binary Arithmetic \& Two's Complement}, arc=3pt, boxrule=1pt]
\footnotesize
\begin{multicols}{2}
\begin{itemize}[noitemsep, leftmargin=*]
  \item Binary $\to$ Decimal: $\sum d_i \times 2^i$
  \item Decimal $\to$ Binary: divide by 2, read remainders bottom-up
  \item Hex $\leftrightarrow$ Binary: group by 4 bits per hex digit
  \item $n$-bit range (signed): $-2^{n-1}$ to $2^{n-1}-1$
\end{itemize}
\columnbreak
\begin{itemize}[noitemsep, leftmargin=*]
  \item \textbf{Two's complement:} invert all bits, then add 1
  \item Subtraction $A - B$ = $A +$ (two's complement of $B$)
  \item Overflow: two positives $\to$ negative (or vice versa)
  \item $1\ \text{byte} = 8\ \text{bits}$; $1\ \text{word (MIPS)} = 4\ \text{bytes} = 32\ \text{bits}$
\end{itemize}
\end{multicols}
\end{tcolorbox}

\vspace{4pt}

% ── Boolean Algebra & Gates ───────────────────────────────────────────────────
\begin{tcolorbox}[colback=part4bg, colframe=sectionorange,
  title={\bfseries Boolean Algebra \& Logic Gates}, arc=3pt, boxrule=1pt]
\footnotesize
\begin{multicols}{2}
\textbf{Key laws:}
\begin{itemize}[noitemsep, leftmargin=*]
  \item Complement: $A\cdot\bar{A}=0$; $A+\bar{A}=1$
  \item De Morgan: $\overline{A\cdot B}=\bar{A}+\bar{B}$; $\overline{A+B}=\bar{A}\cdot\bar{B}$
  \item Absorption: $A+(A\cdot B)=A$
  \item Double negation: $\overline{\overline{A}}=A$
\end{itemize}
\columnbreak
\textbf{Gates (output = 1 when\ldots):}
\begin{itemize}[noitemsep, leftmargin=*]
  \item AND ($A\cdot B$): all inputs are 1
  \item OR ($A+B$): any input is 1
  \item XOR ($A\oplus B$): inputs are \emph{different}
  \item NAND: not (all inputs 1)
  \item NOR: all inputs are 0
  \item Bubble = inversion; NAND/NOR are universal
\end{itemize}
\end{multicols}
\textbf{ANSI:} shapes (D=AND, pointy=OR) \quad \textbf{IEC:} boxes (\& = AND, $\geq$1 = OR)
\end{tcolorbox}

\vspace{4pt}

% ── Circuits ──────────────────────────────────────────────────────────────────
\begin{tcolorbox}[colback=part5bg, colframe=sectionteal,
  title={\bfseries Key Circuits}, arc=3pt, boxrule=1pt]
\footnotesize
\begin{multicols}{2}
\begin{itemize}[noitemsep, leftmargin=*]
  \item \textbf{Half adder:} Sum $= A\oplus B$; Carry $= A\cdot B$
  \item \textbf{Full adder:} adds $A$, $B$, $C_{\mathrm{in}}$; chain for multi-bit
  \item \textbf{MUX:} selects input based on select lines
  \item \textbf{Decoder:} asserts exactly one output line
\end{itemize}
\columnbreak
\begin{itemize}[noitemsep, leftmargin=*]
  \item \textbf{SOP:} OR of AND terms (one per row where out = 1)
  \item \textbf{POS:} AND of OR terms (one per row where out = 0)
  \item \textbf{ALU control} picks operation (add=0010, sub=0110, AND=0000, OR=0001, SLT=0111)
\end{itemize}
\end{multicols}
\end{tcolorbox}

\vspace{4pt}

% ── MIPS ──────────────────────────────────────────────────────────────────────
\begin{tcolorbox}[colback=partbg, colframe=sectionblue,
  title={\bfseries MIPS Assembly Quick Reference}, arc=3pt, boxrule=1pt]
\footnotesize
\begin{multicols}{2}
\textbf{Instruction formats:}
\begin{itemize}[noitemsep, leftmargin=*]
  \item \textbf{R-type:} \texttt{op rd, rs, rt} \quad (result, val1, val2)
  \item \textbf{I-type arith:} \texttt{op rt, rs, imm}
  \item \textbf{I-type mem:} \texttt{lw/sw rt, offset(rs)}\\
        \quad address = \texttt{Md[rs + offset]}
  \item \textbf{J-type:} \texttt{j/jal target}
\end{itemize}
\columnbreak
\textbf{Key registers:}
\begin{itemize}[noitemsep, leftmargin=*]
  \item \texttt{\$zero}: always 0
  \item \texttt{\$a0--\$a3}: function arguments
  \item \texttt{\$v0--\$v1}: return values
  \item \texttt{\$t0--\$t9}: temporaries (not preserved)
  \item \texttt{\$s0--\$s7}: saved (must be preserved)
  \item \texttt{\$sp}: stack pointer (grows down)
  \item \texttt{\$ra}: return address (set by \texttt{jal})
\end{itemize}
\end{multicols}
\divider
\textbf{Memory:} \texttt{Md[addr]} is like \texttt{array[index]}. Words = 4 bytes; must be 4-byte aligned. Array index $i$ is at offset $i \times 4$ from base. \texttt{lw/sw} are the \emph{only} instructions that touch memory.\\[2pt]
\textbf{Call convention:} args in \texttt{\$a0--\$a3}; \texttt{jal} to call; save \texttt{\$s} regs and \texttt{\$ra} on stack if needed; return in \texttt{\$v0}; \texttt{jr \$ra} to return.
\end{tcolorbox}

\newpage

%% =============================================================================
%%  PART VI -- x86-64 Architecture & Assembly
%% =============================================================================
\redsections
\addcontentsline{toc}{section}{\textcolor{sectionred}{PART VI \textemdash\ x86-64 Architecture \& Assembly}}
\addcontentsline{toc}{subsection}{RISC vs.\ CISC: MIPS vs.\ x86-64}
\addcontentsline{toc}{subsection}{x86-64 Register File}
\addcontentsline{toc}{subsection}{System Calls \& Hardware Interfaces}
\addcontentsline{toc}{subsection}{Stack Mechanics \& the Call Frame}
\addcontentsline{toc}{subsection}{Control Flow \& the RFLAGS Register}
\addcontentsline{toc}{subsection}{Memory Security \& Buffer Overflows}
\addcontentsline{toc}{subsection}{x86-64 Instruction Quickstart}
\partbanner{part6bg}{sectionred}{Part VI \quad x86-64 Architecture \& Assembly}

\vspace{6pt}
\begin{center}\small\itshape
  x86-64 is the architecture running on virtually every desktop, laptop, and server.
  Understanding it bridges assembly, the OS kernel, and security.
\end{center}

% -- RISC vs CISC -------------------------------------------------------------
\section*{RISC vs.\ CISC: MIPS vs.\ x86-64}

MIPS is a \textbf{RISC} (Reduced Instruction Set Computer) architecture; x86-64 is \textbf{CISC} (Complex Instruction Set Computer). These design philosophies lead to fundamental differences in every aspect of the architecture.

\noindent
{\renewcommand{\arraystretch}{1.6}
\begin{tabularx}{\textwidth}{>{\bfseries\raggedright\arraybackslash}p{4cm}
                              >{\raggedright\arraybackslash}X
                              >{\raggedright\arraybackslash}X}
\toprule
\rowcolor{tableheaderred}
\textbf{Feature} & \textbf{MIPS (RISC)} & \textbf{x86-64 (CISC)} \\
\midrule
\rowcolor{rowRe} Instruction size      & Fixed: \textbf{exactly 32 bits} always & Variable: \textbf{1--15 bytes} per instruction \\
\rowcolor{rowA}  Instruction count     & Small, uniform set ($\sim$50 core instructions) & Large, irregular set (thousands of encodings) \\
\rowcolor{rowRe} Memory access         & \textbf{Load/store only} --- arithmetic on registers only & Memory operands in arithmetic (\texttt{add [rbx], 1}) \\
\rowcolor{rowA}  Registers             & 32 general-purpose, all equal (\texttt{\$0}--\texttt{\$31}) & 16 general-purpose, historically unequal (\texttt{RAX}--\texttt{R15}) \\
\rowcolor{rowRe} Addressing modes      & Base + offset only & Many: immediate, register, direct, scaled index \\
\rowcolor{rowA}  Pipeline              & Simple, uniform pipeline (designed for it) & Complex, deep out-of-order pipeline \\
\rowcolor{rowRe} Return address        & Stored in \texttt{\$ra} register & Pushed onto the \textbf{stack} by \texttt{CALL} \\
\rowcolor{rowA}  Compiler target       & Common in embedded, routers, gaming consoles & Dominant in PCs, servers, OS kernels \\
\bottomrule
\end{tabularx}}

\begin{tcolorbox}[colback=part6bg, colframe=sectionred, arc=3pt, boxrule=1pt]
\textbf{Why x86-64 matters:} Despite MIPS being cleaner for teaching, x86-64 is what runs your OS, your browser, and virtually every cloud server. When you write C/C++/Rust and it compiles for your laptop, it produces x86-64 machine code. Understanding it is essential for systems programming, debugging at the assembly level, and understanding security vulnerabilities like buffer overflows.
\end{tcolorbox}

% -- x86-64 Register File -----------------------------------------------------
\section*{x86-64 Register File}

x86-64 has \textbf{16 general-purpose registers}, each 64 bits wide. Each register has sub-register aliases for backward compatibility with older code (32-bit, 16-bit, 8-bit). This is a key difference from MIPS which uses clean sequential numbers.

\noindent
{\small\renewcommand{\arraystretch}{1.5}
\begin{tabularx}{\textwidth}{>{\ttfamily\bfseries\centering\arraybackslash}p{1.4cm}
                              >{\ttfamily\centering\arraybackslash}p{1.4cm}
                              >{\ttfamily\centering\arraybackslash}p{1.4cm}
                              >{\ttfamily\centering\arraybackslash}p{1.4cm}
                              >{\raggedright\arraybackslash}X
                              >{\centering\arraybackslash}p{2cm}}
\toprule
\rowcolor{tableheaderred}
\textbf{64-bit} & \textbf{32-bit} & \textbf{16-bit} & \textbf{8-bit} & \textbf{Conventional Use} & \textbf{Call-preserved?} \\
\midrule
\rowcolor{rowRe} RAX & EAX & AX  & AL  & Return value; accumulator for arithmetic      & No \\
\rowcolor{rowA}  RBX & EBX & BX  & BL  & General purpose (``base'' register)            & \textbf{Yes} \\
\rowcolor{rowRe} RCX & ECX & CX  & CL  & 4th argument; counter in loops                 & No \\
\rowcolor{rowA}  RDX & EDX & DX  & DL  & 3rd argument; high bits of \texttt{mul/div}    & No \\
\rowcolor{rowRe} RSI & ESI & SI  & SIL & 2nd argument; source for string ops            & No \\
\rowcolor{rowA}  RDI & EDI & DI  & DIL & 1st argument; destination for string ops       & No \\
\rowcolor{rowRe} RSP & ESP & SP  & SPL & \textbf{Stack pointer} (grows downward)        & \textbf{Yes} \\
\rowcolor{rowA}  RBP & EBP & BP  & BPL & \textbf{Base/frame pointer} (optional in x86-64) & \textbf{Yes} \\
\rowcolor{rowRe} R8  & R8D & R8W & R8B & 5th argument                                   & No \\
\rowcolor{rowA}  R9  & R9D & R9W & R9B & 6th argument                                   & No \\
\rowcolor{rowRe} R10--R11 & \ldots & \ldots & \ldots & Scratch / temporary registers           & No \\
\rowcolor{rowA}  R12--R15 & \ldots & \ldots & \ldots & General purpose                         & \textbf{Yes} \\
\bottomrule
\end{tabularx}}

\vspace{4pt}
\begin{tcolorbox}[colback=part6bg, colframe=sectionred, arc=3pt, boxrule=1pt]
\textbf{Sub-register aliasing:} All four names (\texttt{RAX}/\texttt{EAX}/\texttt{AX}/\texttt{AL}) refer to the \emph{same physical register}, just different widths. Writing to \texttt{EAX} (32-bit) \textbf{zero-extends} into \texttt{RAX}. Writing to \texttt{AX} (16-bit) or \texttt{AL} (8-bit) only modifies those bits --- the upper bits of \texttt{RAX} are unchanged. This is a frequent source of bugs for programmers coming from MIPS.\\[4pt]
\textbf{Calling convention (System V AMD64 -- Linux/macOS):} First 6 integer args go in \texttt{RDI}, \texttt{RSI}, \texttt{RDX}, \texttt{RCX}, \texttt{R8}, \texttt{R9}. Return value in \texttt{RAX}. Compare to MIPS: \texttt{\$a0--\$a3} for 4 args, return in \texttt{\$v0}.
\end{tcolorbox}

% -- System Calls -------------------------------------------------------------
\section*{System Calls \& Hardware Interfaces}

In x86-64 Linux, user programs interact with the OS kernel via \textbf{system calls}. Executing \texttt{syscall} triggers a ring transition from user mode (ring 3) to kernel mode (ring 0), where the kernel handles the request and returns.

\textbf{System call calling convention (Linux x86-64):}
\begin{itemize}[noitemsep]
  \item \textbf{RAX} = system call number (which syscall to invoke)
  \item \textbf{RDI, RSI, RDX, R10, R8, R9} = arguments (in order)
  \item After \texttt{syscall}: \textbf{RAX} holds the return value (negative = error code)
\end{itemize}

\noindent
{\renewcommand{\arraystretch}{1.5}
\begin{tabularx}{\textwidth}{>{\centering\arraybackslash}p{1.5cm}
                              >{\ttfamily\bfseries\raggedright\arraybackslash}p{2.5cm}
                              >{\raggedright\arraybackslash}X
                              >{\raggedright\arraybackslash}p{3.5cm}}
\toprule
\rowcolor{tableheaderred}
\textbf{RAX} & \textbf{Name} & \textbf{Description} & \textbf{Args (RDI, RSI, RDX)} \\
\midrule
\rowcolor{rowRe} 0  & sys\_read   & Read bytes from a file descriptor & fd, buf pointer, count \\
\rowcolor{rowA}  1  & sys\_write  & Write bytes to a file descriptor  & fd, buf pointer, count \\
\rowcolor{rowRe} 2  & sys\_open   & Open a file, returns fd           & filename ptr, flags, mode \\
\rowcolor{rowA}  3  & sys\_close  & Close a file descriptor           & fd \\
\rowcolor{rowRe} 9  & sys\_mmap   & Map memory (allocate heap-like)   & addr, length, prot, flags, fd, offset \\
\rowcolor{rowA}  39 & sys\_getpid & Get current process ID            & (none) \\
\rowcolor{rowRe} 60 & sys\_exit   & Terminate process                 & exit status code \\
\rowcolor{rowA}  231 & sys\_exit\_group & Terminate all threads in group & exit status code \\
\bottomrule
\end{tabularx}}

\begin{lstlisting}[style=x86, caption={Hello World via direct syscall (no C library)}]
section .data
    msg db "Hello, World!", 0x0A  ; string + newline
    len equ $ - msg               ; length = current pos - start of msg

section .text
    global _start

_start:
    ; sys_write(1, msg, len)
    mov rax, 1          ; syscall number: write
    mov rdi, 1          ; arg1: file descriptor 1 = stdout
    mov rsi, msg        ; arg2: pointer to the message
    mov rdx, len        ; arg3: number of bytes to write
    syscall             ; kernel handles the write

    ; sys_exit(0)
    mov rax, 60         ; syscall number: exit
    xor rdi, rdi        ; arg1: exit status = 0 (xor reg,reg is faster than mov reg,0)
    syscall
\end{lstlisting}

% -- Stack Mechanics ----------------------------------------------------------
\section*{Stack Mechanics \& the Call Frame}

In x86-64, the stack \textbf{grows downward} in memory. \texttt{RSP} always points to the top (lowest address) of the current stack frame.

\begin{itemize}[noitemsep]
  \item \textbf{\texttt{PUSH reg}:} Subtracts 8 from \texttt{RSP}, then writes \texttt{reg} to \texttt{[RSP]}. One atomic instruction.
  \item \textbf{\texttt{POP reg}:} Reads \texttt{[RSP]} into \texttt{reg}, then adds 8 to \texttt{RSP}.
  \item \textbf{\texttt{CALL label}:} Pushes the \textbf{return address} (\texttt{RIP+instruction\_size}) onto the stack, then jumps. This is fundamentally different from MIPS \texttt{jal} which stores the return address in \texttt{\$ra}.
  \item \textbf{\texttt{RET}:} Pops the top of the stack into \texttt{RIP} and resumes execution there.
\end{itemize}

\textbf{The standard function prologue and epilogue:}

\begin{lstlisting}[style=x86, caption={Standard x86-64 function frame setup}]
; ---- PROLOGUE: establish the stack frame ----
my_function:
    push rbp            ; save caller's base pointer onto stack
    mov  rbp, rsp       ; set our frame pointer = current stack top
    sub  rsp, 32        ; allocate 32 bytes for local variables
                        ; local var at [rbp - 8], [rbp - 16], etc.

    ; ... function body ...
    mov QWORD PTR [rbp - 8], rdi    ; store first argument as local var

; ---- EPILOGUE: tear down the frame ----
    mov  rsp, rbp       ; restore stack pointer (discard locals)
    pop  rbp            ; restore caller's base pointer
    ret                 ; pop return address into RIP, return to caller
\end{lstlisting}

\begin{tcolorbox}[colback=part6bg, colframe=sectionred, arc=3pt, boxrule=1pt]
\textbf{Why RBP?} RSP moves constantly as you push/pop during a function. Using a fixed \textbf{RBP} as an anchor lets you always find local variables at \textbf{constant negative offsets} from RBP, regardless of how many pushes/pops happen in the function body. Modern compilers often omit RBP with \texttt{-O2} and track locals via fixed RSP offsets instead --- but understanding the RBP-based frame is essential for reading debug builds and stack traces.
\end{tcolorbox}

% -- Control Flow -------------------------------------------------------------
\section*{Control Flow \& the RFLAGS Register}

x86-64 conditional jumps do not compare registers directly (unlike MIPS \texttt{beq}). Instead, the \texttt{CMP} or \texttt{TEST} instruction sets flags in the \textbf{RFLAGS} register, and a subsequent conditional jump reads those flags.

\begin{itemize}[noitemsep]
  \item \textbf{\texttt{CMP a, b}}: Computes \texttt{a - b} and sets flags. Does not store the result.
  \item \textbf{\texttt{TEST a, b}}: Computes \texttt{a AND b} and sets flags. Used to test if a value is zero.
  \item \textbf{Key flags:} ZF (zero flag), SF (sign flag), CF (carry flag), OF (overflow flag).
\end{itemize}

\noindent
{\renewcommand{\arraystretch}{1.5}
\begin{tabularx}{\textwidth}{>{\ttfamily\bfseries\centering\arraybackslash}p{2cm}
                              >{\raggedright\arraybackslash}X
                              >{\raggedright\arraybackslash}p{3.5cm}}
\toprule
\rowcolor{tableheaderred}
\textbf{Jump} & \textbf{Condition} & \textbf{Equivalent C} \\
\midrule
\rowcolor{rowRe} jmp  & Always (unconditional) & \texttt{goto} \\
\rowcolor{rowA}  je / jz & ZF = 1 (result was zero) & \texttt{if (a == b)} \\
\rowcolor{rowRe} jne / jnz & ZF = 0 & \texttt{if (a != b)} \\
\rowcolor{rowA}  jl / jnge & SF $\neq$ OF (signed less than) & \texttt{if (a < b)} \\
\rowcolor{rowRe} jle      & ZF = 1 or SF $\neq$ OF & \texttt{if (a <= b)} \\
\rowcolor{rowA}  jg / jnle & ZF = 0 and SF = OF & \texttt{if (a > b)} \\
\rowcolor{rowRe} jge      & SF = OF & \texttt{if (a >= b)} \\
\bottomrule
\end{tabularx}}

\begin{lstlisting}[style=x86, caption={C control flow translated to x86-64}]
; C: if (rdi > 0) { rax = rdi; } else { rax = -rdi; }
    cmp  rdi, 0         ; set flags: rdi - 0
    jle  .negative      ; jump if rdi <= 0
    mov  rax, rdi       ; positive case
    jmp  .done
.negative:
    neg  rdi            ; negate: rdi = -rdi
    mov  rax, rdi
.done:
    ret

; C: for (int i = 0; i < 10; i++) { sum += i; }
    xor  rax, rax       ; sum = 0  (faster than mov rax, 0)
    xor  rcx, rcx       ; i = 0
.loop:
    cmp  rcx, 10        ; i - 10
    jge  .done          ; if i >= 10, exit
    add  rax, rcx       ; sum += i
    inc  rcx            ; i++
    jmp  .loop
.done:
    ret                 ; return sum in rax
\end{lstlisting}

% -- Memory Security ----------------------------------------------------------
\section*{Memory Security \& Buffer Overflows in x86-64}

The key security implication of x86-64's stack model is that \textbf{the return address lives on the same stack as local variables}. An input buffer that overflows its bounds can overwrite the saved return address.

\textbf{The anatomy of a stack buffer overflow:}

\begin{lstlisting}[style=x86, caption={C code and the dangerous stack layout it creates}]
; C code (compiled to x86-64):
void vulnerable(char *input) {
    char buffer[64];         // 64 bytes on the stack
    strcpy(buffer, input);   // NO bounds check!
}
; Stack layout after the prologue (growing downward):
; [rbp + 8]  = return address  <-- target for an attacker
; [rbp + 0]  = saved RBP
; [rbp - 64] = buffer[0..63]   <-- write starts here
;
; If input is 80 bytes: buffer fills, saved RBP overwritten,
;   return address overwritten with attacker-controlled value.
; When 'ret' executes: RIP = attacker's address --> arbitrary code.
\end{lstlisting}

\textbf{OS-level mitigations} (all discussed in the Security notes):
\begin{itemize}[noitemsep]
  \item \textbf{Stack Canary:} Compiler inserts a random value between the buffer and the return address. Checked before \texttt{ret}; mismatch = program terminates.
  \item \textbf{ASLR:} Stack, heap, and library base addresses randomized on every run. Attacker can't predict where to jump.
  \item \textbf{NX / DEP:} Stack memory is marked non-executable. Injected shellcode in the buffer can't run --- but \textbf{Return-Oriented Programming (ROP)} can still chain existing executable instructions.
  \item \textbf{PIE:} The executable itself loads at a random address, closing the last fixed-address ROP gadget source.
\end{itemize}

\begin{tcolorbox}[colback=part6bg, colframe=sectionred, arc=3pt, boxrule=1pt]
\textbf{Why MIPS doesn't have this specific problem:} MIPS stores the return address in \texttt{\$ra}, a \emph{register}, not on the stack. A stack overflow in MIPS cannot overwrite \texttt{\$ra} directly. However, functions that call other functions \emph{must} save \texttt{\$ra} to the stack (as seen in the factorial example in Part IV) --- and at that point, the same vulnerability exists.
\end{tcolorbox}

% -- x86-64 Instruction Quickstart --------------------------------------------
\section*{x86-64 Instruction Quickstart}

\begin{lstlisting}[style=x86, caption={Most common x86-64 instructions with MIPS comparison}]
; DATA MOVEMENT
mov  rax, rbx       ; rax = rbx            (like: add $t0, $t1, $zero)
mov  rax, 42        ; rax = 42             (like: addi $t0, $zero, 42)
mov  rax, [rbx]     ; rax = Mem[rbx]       (like: lw $t0, 0($t1))
mov  [rbx], rax     ; Mem[rbx] = rax       (like: sw $t0, 0($t1))
mov  rax, [rbx+8]   ; rax = Mem[rbx+8]     (like: lw $t0, 8($t1))
lea  rax, [rbx+rcx*4]  ; rax = rbx + rcx*4  (address calc, no memory read)

; ARITHMETIC
add  rax, rbx       ; rax = rax + rbx
add  rax, 5         ; rax = rax + 5        (like: addi $t0, $t0, 5)
sub  rax, rbx       ; rax = rax - rbx
imul rax, rbx       ; rax = rax * rbx (signed; result in rax)
inc  rax            ; rax++               (like: addi $t0, $t0, 1)
dec  rax            ; rax--
neg  rax            ; rax = -rax

; BITWISE
and  rax, rbx       ; rax = rax AND rbx   (like: and $t0, $t0, $t1)
or   rax, rbx       ; rax = rax OR  rbx
xor  rax, rbx       ; rax = rax XOR rbx
xor  rax, rax       ; IDIOM: zero a register (rax = 0), faster than mov rax,0
not  rax            ; rax = bitwise NOT rax
shl  rax, 2         ; rax = rax << 2      (like: sll $t0, $t0, 2)
shr  rax, 2         ; rax = rax >> 2 (logical, zero-fill)
sar  rax, 2         ; rax = rax >> 2 (arithmetic, sign-fill)

; COMPARISON & BRANCHING
cmp  rax, rbx       ; set flags from rax - rbx (no result stored)
test rax, rax       ; set flags from rax AND rax (fast zero test)
je   label          ; jump if ZF=1 (rax == rbx after cmp)
jne  label          ; jump if ZF=0
jl   label          ; jump if signed less than

; FUNCTION CALLS
call function       ; push RIP+size, jump to function
ret                 ; pop into RIP, return to caller
push rax            ; RSP -= 8; Mem[RSP] = rax
pop  rax            ; rax = Mem[RSP]; RSP += 8
\end{lstlisting}

\newpage

%% =============================================================================
%%  PART VII -- Quick Reference Cheat Sheet (Extended)
%% =============================================================================
\purplesections
\addcontentsline{toc}{section}{\textcolor{sectionpurple}{PART VII \textemdash\ x86-64 Quick Reference}}
\partbanner{part3bg}{sectionpurple}{Part VII \quad x86-64 Quick Reference}

\vspace{4pt}
\begin{center}\small\itshape Quick summary of x86-64 to complement the MIPS cheat sheet in Part V.\end{center}
\vspace{4pt}

\begin{tcolorbox}[colback=part6bg, colframe=sectionred,
  title={\bfseries x86-64 vs.\ MIPS at a Glance}, arc=3pt, boxrule=1pt]
\footnotesize
\begin{multicols}{2}
\textbf{x86-64 key facts:}
\begin{itemize}[noitemsep, leftmargin=*]
  \item 16 registers; 64-bit; sub-register aliases (RAX/EAX/AX/AL)
  \item Variable-length instructions (1--15 bytes) --- CISC
  \item Can use memory directly as operands (\texttt{add [rbx], 1})
  \item Return address pushed onto \textbf{stack} by \texttt{CALL}
  \item Syscall args: RAX=number, RDI, RSI, RDX, R10, R8, R9
  \item Stack grows \textbf{down}; RSP = top; locals at \texttt{[rbp - N]}
\end{itemize}
\columnbreak
\textbf{MIPS vs.\ x86-64:}
\begin{itemize}[noitemsep, leftmargin=*]
  \item \texttt{jal} saves return to \$ra $\leftrightarrow$ \texttt{CALL} pushes to stack
  \item \texttt{beq \$a, \$b, L} $\leftrightarrow$ \texttt{cmp rax, rbx} then \texttt{je L}
  \item \texttt{lw \$t0, 8(\$s0)} $\leftrightarrow$ \texttt{mov rax, [rbx+8]}
  \item \texttt{sw \$t0, 0(\$s0)} $\leftrightarrow$ \texttt{mov [rbx], rax}
  \item \texttt{addi \$t0, \$t0, 5} $\leftrightarrow$ \texttt{add rax, 5}
  \item \texttt{sll \$t0, \$t0, 2} $\leftrightarrow$ \texttt{shl rax, 2}
\end{itemize}
\end{multicols}
\end{tcolorbox}

\vspace{4pt}

\begin{tcolorbox}[colback=part6bg, colframe=sectionred,
  title={\bfseries x86-64 Stack \& Memory Security}, arc=3pt, boxrule=1pt]
\footnotesize
\begin{multicols}{2}
\begin{itemize}[noitemsep, leftmargin=*]
  \item \textbf{Buffer overflow}: overflow local var $\to$ overwrite return address on stack
  \item \textbf{CALL}: pushes RIP+N; \textbf{RET}: pops RIP. Return addr is on the stack!
  \item \textbf{Stack canary}: random value between buffer and return addr
  \item \textbf{ASLR}: randomizes stack/heap/lib base addresses
\end{itemize}
\columnbreak
\begin{itemize}[noitemsep, leftmargin=*]
  \item \textbf{NX/DEP}: stack not executable (blocks shellcode injection)
  \item \textbf{ROP}: reuses existing executable code gadgets -- bypasses NX
  \item \textbf{MIPS difference}: \$ra is a register, not on stack (unless callee saves it)
  \item \textbf{Safe practice}: always use \texttt{strncpy}/\texttt{snprintf}; never \texttt{strcpy}/\texttt{gets}
\end{itemize}
\end{multicols}
\end{tcolorbox}

\vspace{4pt}

\begin{tcolorbox}[colback=part6bg, colframe=sectionred,
  title={\bfseries Common x86-64 Idioms}, arc=3pt, boxrule=1pt]
\footnotesize
\begin{multicols}{2}
\begin{itemize}[noitemsep, leftmargin=*]
  \item Zero a register: \texttt{xor rax, rax} (faster than \texttt{mov rax, 0})
  \item Test for zero: \texttt{test rax, rax} then \texttt{jz}
  \item Copy register: \texttt{mov rax, rbx}
  \item \texttt{lea rax, [rbx + rcx*8]}: address math without memory read
  \item \texttt{neg rax}: negate (two's complement)
\end{itemize}
\columnbreak
\begin{itemize}[noitemsep, leftmargin=*]
  \item Prologue: \texttt{push rbp; mov rbp, rsp; sub rsp, N}
  \item Epilogue: \texttt{mov rsp, rbp; pop rbp; ret}
  \item Syscall write stdout: \texttt{mov rax,1; mov rdi,1; syscall}
  \item Syscall exit: \texttt{mov rax,60; xor rdi,rdi; syscall}
  \item \texttt{sar rax, 63}: fill all bits with sign bit (all 0s or all 1s)
\end{itemize}
\end{multicols}
\end{tcolorbox}

\end{document}